\section{Introduction}
\label{sec:introduction}

Forecasting decisions require more than a point estimate when the costs of under- and over-prediction differ. Predictive distributions express uncertainty across the forecast horizon and support decisions that depend on tail risk, quantiles, or intervals. Modern probabilistic forecasting therefore evaluates the entire distribution, not only its center \citep{gneiting2007proper,gneiting2014probabilistic,salinas2020deepar}.

Historical observations are often incomplete. Missing values can remove recent evidence, entire channels, or contiguous blocks, changing both point errors and the residual shape available to a probability model. Recurrent imputation methods, deep generative imputers, and masked forecasting architectures address different parts of this problem \citep{cao2018brits,che2018grud,fortuin2020gpvae,tashiro2021csdi,du2023saits,lee2025timeperceiver}, but Gaussian residual assumptions can remain fragile when errors become heavy-tailed.

Probabilistic forecasts should be compared with proper scores that reward honest distributions. Negative log-likelihood (NLL) is distribution-sensitive, while the continuous ranked probability score (CRPS) evaluates calibration and sharpness together \citep{matheson1976crps,gneiting2007proper,jordan2019scoringrules}. Mean squared error (MSE) and mean absolute error (MAE) remain useful point criteria, but neither replaces distributional evaluation. Coverage and width are likewise diagnostics rather than complete quality measures in isolation.

Paired comparisons also require identical evaluation inputs. If Gaussian and heavy-tailed heads are evaluated under independently sampled masks, mask variation can be mistaken for a distribution-head effect. We therefore fix train, validation, and test splits; fit normalization on training data only; and reuse the same validation and test masks within every comparison.

The proposed final specification is deliberately small: a fixed \StudentT{} predictive marginal with $\df=8$. Gaussian, Laplace, and fixed \StudentT{} variants have identical trainable location and scale projections; the fixed head adds no trainable degrees-of-freedom parameter. The value $8$ is a reference setting that is empirically sufficient in the distribution-head ablation, not a theoretically optimal or dataset-specific choice.

The evidence combines missing-history degradation, four probability-head seeds on fixed TimePerceiver checkpoints, a distribution-family ablation, a PatchTST replication, efficiency measurements, and retained negative results. Our contributions are: (i) a fixed-mask protocol for paired probabilistic forecasting under missing historical observations; (ii) controlled evidence that a matched-capacity fixed heavy-tailed head improves proper scores on the evaluated TimePerceiver conditions; and (iii) a transparent second-backbone assessment showing that probabilistic-score gains transfer while a small Electricity point-error tradeoff prevents a complete multi-metric conclusion.
